\section{Fekete limits for the actual test-configuration base ideals}

\begin{lemma}\label{lem:model-ideal-norm-monotone}
\leanok
\lean{EnvelopeProperty.Berkovich.RelativeSchemeValuedPoint.modelIdealNorm_mono}
\lean{EnvelopeProperty.Berkovich.RelativeSchemeValuedPoint.modelIdealLog_mono}
For a proper scheme over a bounded valued affine base, inclusion of actual
ideal sheaves implies the same inequality between their actual model norms
at every original relative valued point. Their logarithms obey the same
monotonicity.
\end{lemma}
\begin{proof}
Evaluate both ideals on any affine neighborhood of the actual valuative
center. Original ideal-sheaf inclusion implies inclusion of the affine
ideals. Their norms are monotone because they are suprema of the values
of their elements. The intrinsic model norms equal these affine norms.
Apply the monotonicity of the extended-real logarithm for the last assertion.
\end{proof}

\begin{lemma}\label{lem:native-power-base-ideal-supermultiplicative}
\leanok
\lean{EnvelopeProperty.Berkovich.RelativeSchemeValuedPoint.powerBaseIdeal_norm_mul_le}
Let $M$ be an actual invertible sheaf on a proper model. For the base ideals
$I_m$ of all global sections of the original powers $M^m$,
$\|I_m\|\|I_n\|\leq\|I_{m+n}\|$ at every actual point, including zero norms.
\uses{lem:model-ideal-norm-monotone,prop:actual-line-power-base-ideal-grading}
\end{lemma}
\begin{proof}
Actual tensor multiplication gives the whole original ideal inclusion
$I_m I_n\subset I_{m+n}$. Intrinsic model ideal norms multiply on actual
ideal products. The preceding lemma gives the claimed inequality.
\end{proof}

\begin{theorem}\label{thm:native-test-configuration-fekete-limit}
\leanok
\lean{EnvelopeProperty.Berkovich.RelativeSchemeValuedPoint.powerBaseIdealNegativeLog}
\lean{EnvelopeProperty.NativeTestConfiguration.TestConfiguration.exists_powerBaseIdeal_fekete_limit}
Take the original ample fibre line and original source line-action and
fibre-one data of eventual verticality, and fix any bounded parameter
valuation with nonzero parameter value. At every original native point,
the actual sequence $-\log(\|I_m\|)/m$ has a finite nonnegative limit.
After a single bound derived from the original ampleness, that limit is
precisely the infimum of these normalized negative logarithms on the tail.
No positivity, norm, subadditivity, boundedness or limit certificate is
supplied as an additional input.
\uses{thm:native-test-configuration-eventual-positive-base-ideal-norm,lem:native-power-base-ideal-supermultiplicative,thm:tail-subadditive-fekete}
\end{theorem}
\begin{proof}
Original ampleness and source data give one bound $N$ after which every
$I_m$ has strictly positive norm at every native point. All model ideal
norms are at most one. Thus their real negative logarithms are nonnegative,
and the preceding supermultiplicativity inequality becomes subadditivity
for $m,n\geq N$ after taking real logarithms and changing sign. The sum
$m+n$ also belongs to the same tail, so all logarithms in that inequality
are finite. Apply the tail version of Fekete convergence, which needs no
condition on the possibly vanishing finitely many earlier base ideals.
The early real-logarithm convention at zero does not affect the limit.
This is the actual ideal-norm limit; comparison with the original weighted
Laurent-section Fubini--Study envelope is a separate obligation.
\end{proof}


\begin{corollary}\label{cor:native-test-configuration-root-limit}
\leanok
\lean{EnvelopeProperty.NativeTestConfiguration.TestConfiguration.exists_powerBaseIdeal_root_limit}
For each actual native point, the real roots $\|I_m\|^{1/m}$ converge to
some $b$ with $0<b\leq1$.
\uses{thm:native-test-configuration-fekete-limit}
\end{corollary}
\begin{proof}
Let $a\geq0$ be the actual negative-logarithm limit from the theorem.
Put $b=\exp(-a)$, which is positive and at most one. On the derived
positive tail, the actual roots equal the exponential of the negative
normalized negative logarithms. Exponential continuity and the preceding
limit imply convergence; the finitely many earlier roots have no effect.
\end{proof}
